% Rozdzial 21. Samodzielna tresc rozdzialu do \input (bez preambuly).
% Wymaga listings, enumitem, tabularx, booktabs, hyperref, makr
% \code, \term, \nota. Przyklady: standard oznaczony przy kazdym kodzie.
% Stan omowienia: pazdziernik 2026. Nalezy sprawdzac stan wdrozenia API.
\chapter{Nowości C++23 i C++26: język, biblioteka i przenośność}
\label{chap:cpp21-new}

\section{Jak czytać informacje o nowym standardzie?}

Zatwierdzenie funkcjonalności przez komitet C++ nie oznacza, że
wszystkie kompilatory i biblioteki standardowe już ją dostarczają.
W szczególności trzeba odróżniać: (1) zgodność kompilatora z regułami
\emph{języka}, (2) implementację odpowiedniego nagłówka w \emph{bibliotece
standardowej} oraz (3) gotowość systemu budowania i pakietowania.
Opisy w tym rozdziale dotyczą C++23 i pakietu funkcji C++26
według stanu na październik 2026; dostępność należy sprawdzić dla
konkretnego środowiska GCC/libstdc++, Clang/libc++ lub MSVC/STL.

\begin{center}\small
\begin{tabularx}{\textwidth}{@{}p{0.30\textwidth}X@{}}
\toprule
Warstwa & Praktyczna weryfikacja \\
\midrule
Standard języka & \code{-std=c++23}, \code{-std=c++2c} (lub wariant kompilatora), \code{__cplusplus} i makra \code{__cpp_*}. \\
Biblioteka standardowa & \code{<version>}, makra \code{__cpp_lib_*}, dostępność nagłówków i linkowania. \\
Narzędzia budowania & Wsparcie CMake, modułów, diagnostyki i zależności BMI. \\
Platforma & Implementacja wielowątkowości, wektoryzacji, I/O, architektury procesora. \\
\bottomrule
\end{tabularx}
\end{center}

\subsection{Feature-test macros zamiast testowania numeru kompilatora}

Nie należy zakładać, że \code{__cplusplus >= 202302L} oznacza obecność
każdego nowego nagłówka. Wykrywamy \emph{konkretną funkcję}.
\code{__has_include} sprawdza wyłącznie obecność pliku nagłówkowego;
nie dowodzi, że zawartość została w pełni zaimplementowana.

\begin{lstlisting}[caption={Bezpieczny wybor std::expected albo wariantu zgodnego z C++20.}]
#include <version>
#include <iostream>

#if defined(__cpp_lib_expected) && __cpp_lib_expected >= 202202L
    #include <expected>
    #define HAS_EXPECTED 1
#else
    #include <optional>
    #define HAS_EXPECTED 0
#endif

int main() {
#if HAS_EXPECTED
    std::expected<int, int> result = 42;
    std::cout << result.value() << '\n';
#else
    std::optional<int> result = 42;
    std::cout << *result << '\n';
#endif
}
\end{lstlisting}

Wartości makr zmieniają się, kiedy standard rozszerza dane API.
W bibliotekach produkcyjnych zaleca się test konfiguracji
\emph{compile-and-link}, a nie tylko sprawdzenie preprocesora.

\section{C++23: \texttt{std::expected<T,E>} i jawne błędy}

\begin{center}\small
\begin{tabularx}{\textwidth}{@{}p{0.35\textwidth}X@{}}
\toprule
Operacja & Znaczenie \\
\midrule
\code{std::expected<T,E>} & Przechowuje wartość \code{T} albo błąd \code{E}. \\
\code{std::unexpected(error)} & Tworzy wariant błędu przekazywany do \code{expected}. \\
\code{has_value()} & Sprawdza, czy dostępna jest wartość, a nie błąd. \\
\code{value()} & Zwraca wartość; gdy przechowywany jest błąd, zgłasza \code{std::bad_expected_access}. \\
\code{error()} & Udostępnia błąd; należy używać, gdy obiekt przechowuje błąd. \\
\code{value_or(fallback)} & Zwraca wartość albo podany zamiennik. \\
\bottomrule
\end{tabularx}
\end{center}
\label{sec:cpp21-expected}

\code{std::expected<T,E>} reprezentuje wynik typu \code{T} albo błąd
typu \code{E}. W przeciwieństwie do \code{std::optional<T>} przenosi
informację, \emph{dlaczego} operacja się nie powiodła. Nie jest
zamiennikiem wyjątków w każdej sytuacji: użytkownik musi obsłużyć
stan błędu, a \code{value()} może zgłosić \code{bad_expected_access}.

\begin{lstlisting}[caption={Parser z jawnym typem bledu (kompletny program C++23).}]
#include <charconv>
#include <expected>
#include <iostream>
#include <string>
#include <string_view>

std::expected<int, std::string> parse_integer(std::string_view text) {
    int value{};
    auto [ptr, error] = std::from_chars(text.data(),
                                        text.data() + text.size(),
                                        value);
    if (error != std::errc{} || ptr != text.data() + text.size())
        return std::unexpected("invalid integer");
    return value;
}

int main() {
    auto result = parse_integer("123");
    if (result) std::cout << *result << '\n';
    else std::cerr << result.error() << '\n';
}
\end{lstlisting}

Nowsze operacje monadyczne \code{and_then}, \code{transform},
\code{or_else} i \code{transform_error} można wykorzystywać po
sprawdzeniu właściwego poziomu makra \code{__cpp_lib_expected}.
\code{expected} daje się również stosować z \code{void} jako typem
sukcesu, co jest przydatne w operacjach zapisujących dane.

\section{C++23: \texttt{std::print}, \texttt{std::format} i formatowanie}

\code{std::format} jest częścią C++20, natomiast \code{std::print}
i \code{std::println} wprowadzono w C++23. Ich wsparcie w praktyce
zależy głównie od wersji biblioteki standardowej. Formatowanie
stabilnych protokołów danych nadal wymaga jawnego kontraktu, a nie
przypadkowego wydruku do terminala.

\begin{lstlisting}[caption={Formatowany wydruk C++23 (wymaga naglowka <print>).}]
#include <print>

int main() {
    std::println("value = {:.2f}", 3.14159);
    std::print("without newline: {}", 42);
}
\end{lstlisting}

\code{<format>} oferuje bogatszy język formatowania niż ręczne
sklejanie napisów, ale jego użycie dla danych nieufnych wymaga
kontroli formatu i argumentów. Należy odróżniać tekst formatu
znany przy kompilacji od wariantów wykorzystujących format dynamiczny.

\section{C++23: \texttt{std::mdspan} i macierze wielowymiarowe}
\label{sec:cpp21-mdspan}

\code{std::mdspan} jest niewłasnościowym widokiem na pamięć
wielowymiarową. Podobnie jak \code{std::span} nie zarządza czasem życia
bufora. Pozwala opisać ekstensje, odwzorowanie indeksów na adresy
i sposób dostępu. Układy \code{layout_right} i \code{layout_left}
mają inny sposób wyliczania przesunięcia w pamięci; nie należy
przypisywać im automatycznie takiej samej wydajności.

\begin{lstlisting}[caption={Macierz 2 x 3 przechowywana w vector (C++23).}]
#include <mdspan>
#include <vector>
#include <iostream>

int main() {
    std::vector<double> a{1, 2, 3, 4, 5, 6};
    std::mdspan matrix(a.data(), 2, 3);
    matrix[1, 2] = 10.0;  // operator indeksowania C++23
    std::cout << a[5] << '\n';
}
\end{lstlisting}

\nota{Dokładna składnia operatora indeksowania i dostępność nagłówka
\code{<mdspan>} mogą być zależne od implementacji i późniejszych
poprawek. Nie wolno zwracać \code{mdspan} do lokalnego wektora,
który zostanie zniszczony po wyjściu z funkcji.}

\section{C++23: generatory, wymazywanie typu i inne klasy użytkowe}

\subsection{\texttt{std::generator}: od korutyny do zakresu}

\code{std::generator<T>} z nagłówka \code{<generator>} jest standardowym
narzędziem C++23 do leniwego przekazywania wartości przez \code{co_yield}.
Nie każdy kompilator udostępnia tę bibliotekę nawet wtedy, gdy
obsługuje składnię korutyn C++20. Generator nie tworzy automatycznie
oddzielnego wątku ani mechanizmu planowania asynchronicznego.

\begin{lstlisting}[caption={Leniwy generator liczb calkowitych (C++23; wymaga <generator>).}]
#include <generator>
#include <iostream>

std::generator<int> numbers(int limit) {
    for (int i = 0; i < limit; ++i)
        co_yield i;
}

int main() {
    for (int v : numbers(4))
        std::cout << v << ' ';
}
\end{lstlisting}

\subsection{\texttt{move\_only\_function} i koszt własności callbacku}

\code{std::move_only_function<R(Args...)>} jest odpowiednikiem
opisu wywoływalnego obiektu, który \emph{nie wymaga kopiowalności}
przechwyconego stanu. Ma zastosowanie m.in. w kolejkach zadań
przechowujących lambdy zawierające \code{unique_ptr}.

\begin{lstlisting}[caption={Callback z obiektem posiadajacym wylaczna wlasnosc (C++23).}]
#include <functional>
#include <iostream>
#include <memory>
#include <utility>

int main() {
    auto payload = std::make_unique<int>(42);
    std::move_only_function<int()> get =
        [p = std::move(payload)] { return *p; };
    std::cout << get() << '\n';
}
\end{lstlisting}

Ponadto C++23 dodaje m.in. \code{std::byteswap}, \code{std::to_underlying},
\code{std::forward_like}, \code{std::bind_back}, \code{std::flat_map},
\code{std::flat_set}, \code{std::stacktrace} i \code{<spanstream>}.
Każda z tych konstrukcji ma własny zakres użycia i niezależne
wymagania implementacyjne.

\section{C++23: Ranges i materializacja widoków}
\label{sec:cpp21-ranges}

C++23 rozbudowuje bibliotekę Ranges m.in. o \code{views::zip},
\code{views::chunk}, \code{views::slide}, \code{views::stride},
\code{views::enumerate} i \code{ranges::to}. Część algorytmów
czyta wyniki leniwie; \code{ranges::to} materializuje widok w
kontenerze. To duża zmiana ergonomii, ale reguły czasu życia
widoków i koszt materializacji nadal obowiązują.

\begin{lstlisting}[caption={Przekazanie wynikow widoku do wektora (C++23).}]
#include <iostream>
#include <ranges>
#include <vector>

int main() {
    auto values = std::views::iota(0, 10)
                | std::views::filter([](int x) { return x % 2 == 0; });
    auto result = std::ranges::to<std::vector<int>>(values);
    std::cout << result.size() << '\n'; // 5
}
\end{lstlisting}

Dla funkcji, których implementacja nie wspiera jeszcze \code{ranges::to},
można zastosować ręczną iterację i \code{push_back} albo
\code{ranges::copy} wraz z odpowiednim inserterem.

\section{C++23: składnia języka --- deducing this i \texttt{if consteval}}

\emph{Explicit object parameter} (często nazywany \emph{deducing this})
pozwala jawnie opisać obiekt, na którym wywoływana jest funkcja
składowa. Może zastąpić rodzinę przeciążeń kwalifikowanych
\code{&}, \code{const&} oraz \code{&&}, ale nadmierna ogólność
może uczynić kontrakt funkcji mniej czytelnym.

\begin{lstlisting}[caption={Jawny parametr obiektu (C++23).}]
#include <iostream>

struct Counter {
    int value = 0;
    void add(this Counter& self, int step) {
        self.value += step;
    }
};

int main() {
    Counter x;
    x.add(3);
    std::cout << x.value << '\n';
}
\end{lstlisting}

\code{if consteval} pozwala wykonać inną gałąź podczas ewaluacji
stałej, bez traktowania zwykłej funkcji czasu wykonywania jako
\code{consteval}. Trzeba jednak nadal pamiętać o ograniczeniach
konstruowania stałych wyrażeń.

\begin{lstlisting}[caption={if consteval (C++23, kompletny program).}]
#include <iostream>

constexpr int double_value(int n) {
    if consteval {
        return n * 2;
    } else {
        return n + n;
    }
}
static_assert(double_value(21) == 42);
int main() { std::cout << double_value(3) << '\n'; }
\end{lstlisting}

W rzeczywistym projekcie rozdzielenie ścieżek powinno być
uzasadnione innym dostępnym zestawem mechanizmów na etapie
kompilacji i wykonania, a nie sztucznym powielaniem identycznego kodu.

\section{C++26: stały rozmiar przechowywania bez alokacji}
\label{sec:cpp21-inplace}

\code{std::inplace_vector<T,N>} jest kontenerem o zmiennym rozmiarze,
którego \emph{maksymalna pojemność} wynosi \code{N}, a pamięć dla
elementów jest umieszczona w samym obiekcie. Różni się więc od
\code{std::array<T,N>}, która zawsze zawiera \code{N} elementów.
Nie jest zwykłym \code{vector}: przy przekroczeniu pojemności
nie może zwiększyć jej przez alokację.

\begin{lstlisting}[caption={Kontener inplace\_vector (C++26; wymagane wsparcie biblioteki).}]
#include <inplace_vector>
#include <iostream>

int main() {
    std::inplace_vector<int, 8> values;
    values.push_back(10);
    values.push_back(20);
    std::cout << values.size() << '\n';
}
\end{lstlisting}

\code{std::hive} to kolejna nowa struktura kontenerowa C++26,
ukierunkowana na efektywne wstawianie i usuwanie elementów oraz
określone gwarancje stabilności istniejących elementów.
Nie stanowi zamiennika \code{vector} dla obliczeń wymagających
ciągłego bloku pamięci i lokalności odczytów.

\section{C++26: SIMD, algebra liniowa i wielowątkowość}
\label{sec:cpp21-simd}

Nagłówek \code{<simd>} i typy \code{std::simd} służą do jawnego
wyrażania obliczeń na paczkach wartości, potencjalnie mapowanych
na instrukcje wektorowe procesora. Ich wykorzystanie nie zwalnia
z pomiarów; wydajność zależy od architektury, wyrównania,
przepustowości pamięci i zachowania kompilatora.

\begin{lstlisting}[caption={Schemat obliczenia na SIMD (C++26, fragment pogladowy).}]
// #include <simd>
// Wymaga biblioteki z implementacja C++26.
// std::simd<float> a(1.0f), b(2.0f);
// auto c = a + b;
// Typ i sposob ladowania zalezy od przyjetego ABI/implementacji.
\end{lstlisting}

C++26 rozwija także narzędzia algebry liniowej w \code{<linalg>}
i mechanizmy asynchroniczne oparte na \emph{senders/receivers}
(\code{<execution>}). Nowy model wykonania opisuje kompozycję zadań,
obsługę błędów, anulowanie i relację z harmonogramem pracy.
Sama obecność nagłówka nie dowodzi, że dany executor lub scheduler
działa w pożądany sposób; szczegółowe semantyki zależą od wybranych
operacji i implementacji.

\subsection{Bezpieczne odzyskiwanie pamięci}

Nagłówki \code{<hazard_pointer>} i \code{<rcu>} w C++26 mają pomóc
w zarządzaniu czasem życia obiektów w strukturach współbieżnych.
Nie usuwają one automatycznie problemu ABA ani nie sprawiają,
że dowolna struktura z \code{std::atomic<T*>} jest bezpieczna.
Należy osobno uzasadnić porządek pamięci, gwarancje postępu
i moment fizycznego zniszczenia obiektów (por. rozdział 16).

\section{C++26: język -- pack indexing, refleksja i kontrakty}

C++26 wprowadza mechanizm \emph{pack indexing}, czyli wybór
określonego elementu pakietu parametrów bez rekurencyjnego
rozpakowywania. Jego wsparcie zależy od kompilatora.

\begin{lstlisting}[caption={Indeksowanie pakietu typow (C++26, tylko kompilator ze wsparciem).}]
template<class... Ts>
using first_type = Ts...[0];

static_assert(sizeof(first_type<int, double>) == sizeof(int));
\end{lstlisting}

W C++26 rozwijane są również: \emph{static reflection} (nagłówek
\code{<meta>} i odpowiednia składnia refleksji) oraz
\emph{kontrakty} (np. asercje kontraktowe i wsparcie z
\code{<contracts>}). Są to znaczące nowe możliwości,
ale składnia i zakres praktycznego użycia wymagają konkretnego
kompilatora z odpowiednim wsparciem. Nie podajemy tu pozornie
przenośnego przykładu refleksji lub kontraktów, który mógłby
nie kompilować się w typowej instalacji C++26.

\subsection{Co dzieje się z \texttt{constexpr}?}

W kolejnych rewizjach standardu zwiększa się zakres dozwolonych
obliczeń stałych. Nie wynika z tego, że wszystkie operacje I/O,
zasoby systemowe i dynamiczne struktury mogą automatycznie istnieć
w stałym wyrażeniu. Warto oddzielać możliwość \emph{wykonania obliczenia}
przy kompilacji od zdolności utworzenia \emph{trwałego obiektu}
z zasobami zarządzanymi dynamicznie.

\section{Jak wdrażać C++23/26 do istniejącej biblioteki?}
\label{sec:cpp21-migration}

Włączenie nowego standardu dla całego projektu z dnia na dzień jest
często mniej korzystne niż migracja etapami. Proponujemy strategię:

\begin{enumerate}
\item Spisać minimalne wersje obsługiwanych GCC, Clang, MSVC oraz
      implementacji bibliotek standardowych.
\item Dla każdego nowego API dodać test \emph{compile-and-link},
      a nie tylko sprawdzenie \code{__cplusplus}.
\item Rozróżnić funkcje opcjonalne (np. wygodniejszy wydruk) i wymagane
      dla poprawności (np. określony model zarządzania pamięcią).
\item Zachować wspierany wariant alternatywny, jeśli API ma nadal działać
      z C++20 lub starszym środowiskiem.
\item Wprowadzić osobne zadania CI z różnymi kompilatorami,
      architekturami i wariantami standardowej biblioteki.
\item Przed zmianą interfejsu publicznego przeanalizować kompatybilność
      źródłową, ABI, wyjątki, alokację i obietnice dotyczące czasu życia.
\end{enumerate}

\begin{lstlisting}[caption={Weryfikacja obecnosci nowych funkcji biblioteki (C++23).}]
#include <iostream>
#include <version>

int main() {
#ifdef __cpp_lib_expected
    std::cout << "expected: " << __cpp_lib_expected << '\n';
#else
    std::cout << "expected unavailable\n";
#endif
#ifdef __cpp_lib_mdspan
    std::cout << "mdspan: " << __cpp_lib_mdspan << '\n';
#else
    std::cout << "mdspan unavailable\n";
#endif
#ifdef __cpp_lib_print
    std::cout << "print: " << __cpp_lib_print << '\n';
#else
    std::cout << "print unavailable\n";
#endif
}
\end{lstlisting}

\section{Zestawienie funkcjonalności i rekomendacje}
\begin{center}\footnotesize
\begin{tabularx}{\textwidth}{@{}p{0.33\textwidth}p{0.11\textwidth}X@{}}
\toprule
Funkcjonalność & Standard & Główne zastosowanie \\
\midrule
\code{expected}, \code{print} & C++23 & Jawny wynik/błąd; formatowane I/O. \\
\code{mdspan}, \code{generator} & C++23 & Widoki wielowymiarowe; leniwe zakresy korutyn. \\
\code{move_only_function} & C++23 & Wywoływalne obiekty tylko przenoszalne. \\
\code{ranges::to}, \code{views::zip} & C++23 & Składanie i materializacja zakresów. \\
Deducing this, \code{if consteval} & C++23 & Ergonomia szablonów i ewaluacja stała. \\
\code{inplace_vector}, \code{hive} & C++26 & Kontenery o szczególnych gwarancjach pamięci i stabilności. \\
SIMD, linalg, senders/receivers & C++26 & Obliczenia i asynchroniczne wykonanie. \\
Refleksja, kontrakty & C++26 & Metaprogramowanie i opis warunków poprawności. \\
\bottomrule
\end{tabularx}
\end{center}

\nota{Powyższe przyporządkowanie do wersji standardu nie jest
obietnicą dostępności w dowolnej instalacji C++23/C++26. Najlepszym
źródłem danych o narzędziu jest wynik własnego testu kompilacji,
zgodny z dokumentacją danej standardowej biblioteki.}

\section{C++23: kontenery uporządkowane o zwartej reprezentacji}

\code{std::flat_map} i \code{std::flat_set} to adaptery kontenerowe,
których klucze i wartości są przechowywane w sekwencyjnych
strukturach danych, a nie w węzłach drzewa takich jak \code{std::map}.
Pozwala to uzyskać dobrą lokalność danych i efektywne wyszukiwanie,
zwłaszcza gdy zbiór jest stosunkowo statyczny. Z drugiej strony
wstawianie lub usuwanie w środku może wymagać przesuwania elementów.

\begin{lstlisting}[caption={Mapa o sekwencyjnym przechowywaniu (C++23).}]
#include <flat_map>
#include <iostream>

int main() {
    std::flat_map<int, int> lookup;
    lookup.emplace(2, 20);
    lookup.emplace(1, 10);
    std::cout << lookup.at(2) << '\n';
}
\end{lstlisting}

Dla zastosowań numerycznych często opłaca się samemu porównać
\code{flat_map}, \code{unordered_map} i posortowany \code{vector}
z wyszukiwaniem \code{lower_bound}. Wynik zależy m.in. od liczby
rekordów, częstotliwości modyfikacji i wielkości elementów.

\section{C++23: \texttt{spanstream}, \texttt{stacktrace} i funkcje pomocnicze}

\code{std::spanstream} umożliwia operacje strumieniowe na zewnętrznym,
ciągłym buforze. Nie przejmuje jego własności, dlatego po zakończeniu
czasu życia bufora korzystanie ze strumienia jest błędem.
\code{std::stacktrace} ułatwia uzyskanie śladu wywołań; jakość
symbolizacji i obecność informacji o nazwach funkcji są zależne
od platformy, opcji kompilacji oraz sposobu linkowania.

Mniejsze, lecz praktyczne ułatwienia obejmują:
\begin{description}[style=nextline]
\item[\code{std::byteswap}] Odwrócenie kolejności bajtów dla
odpowiedniego typu całkowitego, przydatne w kodowaniu protokołów.
\item[\code{std::to_underlying}] Jawna konwersja wartości typu wyliczeniowego
na jego typ bazowy.
\item[\code{std::forward_like}] Przenoszenie wzorca kwalifikatorów
\code{const} i kategorii wartości z jednego typu na inne wyrażenie,
przydatne przy pisaniu generycznych adapterów.
\item[\code{std::bind_back}] Wygodne częściowe wiązanie argumentów
od prawej strony listy parametrów.
\end{description}

\begin{lstlisting}[caption={Konwersja enum i zamiana kolejnosci bajtow (C++23).}]
#include <bit>
#include <cstdint>
#include <iostream>
#include <utility>

enum class Status : std::uint32_t { ready = 0x12345678u };

int main() {
    auto raw = std::to_underlying(Status::ready);
    auto swapped = std::byteswap(raw);
    std::cout << std::hex << raw << ' ' << swapped << '\n';
}
\end{lstlisting}

\section{C++26: \texttt{function\_ref} i funkcje bez własności}

\code{std::function_ref} reprezentuje niewłasnościowy widok na
obiekt wywoływalny. Przydaje się, gdy funkcja przyjmuje callback
na czas jednego wywołania i nie zamierza go zachować.
W odróżnieniu od \code{std::function}, nie przejmuje własności
lambdy ani jej przechwyconych danych. Podstawowym zagrożeniem
jest zatem pozostawienie widoku po zakończeniu życia callbacku.

\begin{lstlisting}[caption={Przyklad zastosowania function\_ref (C++26; fragment).}]
#include <functional>

void repeat(int n, std::function_ref<void(int)> callback) {
    for (int i = 0; i < n; ++i)
        callback(i);
}
\end{lstlisting}

Jest to typ zgodny z celem opisanym w rozdziale 20: odróżniamy
przekazanie callbacku na czas wywołania od przechowywania go we
współdzielonym stanie zadań.

\section{C++26: interpretacja nowego modelu wykonania}

Model \emph{senders/receivers} wprowadza pojęcia opisujące wynik
operacji asynchronicznej oraz jej kontynuację. W odróżnieniu od
prostego \code{std::async}, wyrażenie złożone z wielu etapów może
przenosić informacje o błędzie i anulowaniu, a wykonanie może być
zależne od schedulera dostarczonego przez bibliotekę.

Nie należy jednak utożsamiać wyrażenia opisującego zadanie z natychmiast
uruchomioną pracą. Ważne są trzy oddzielne momenty:
\begin{enumerate}
\item utworzenie opisu obliczenia (\emph{sender});
\item połączenie z odbiorcą wyniku (\emph{receiver});
\item rozpoczęcie operacji w wybranym środowisku wykonawczym.
\end{enumerate}

Gdy aplikacja ma już własny executor i korutyny, migracja powinna
zacząć się od próbnego adaptera i testów anulowania, a nie od
mechanicznego zastąpienia wszystkich \code{future}. Biblioteki
systemowe mogą udostępniać własne executory, często o innym modelu
kolejek i innych gwarancjach zamknięcia.

\section{Eksperymentalne wdrażanie nowego API bez utraty przenośności}

Niektóre biblioteki \code{<experimental/...>} dostarczały wybrane
mechanizmy przed przyjęciem ich do standardu. Nie należy zakładać,
że ich typy, przestrzenie nazw czy ABI są identyczne ze standardowymi.
Dla kodu produkcyjnego możliwe są trzy strategie:

\begin{enumerate}
\item Pozostać przy C++20, używając bibliotek zewnętrznych dla brakujących
funkcji i oddzielając te zależności od interfejsu publicznego.
\item Wymagać C++23 i wskazanych minimalnych wersji standardowej
biblioteki, walidowanych w systemie CI.
\item Udostępniać implementację warunkową dla różnych standardów,
przy czym wszystkie warianty muszą zachowywać ten sam kontrakt i
przechodzić wspólny zestaw testów.
\end{enumerate}

Dla funkcji odwołujących się do kompilatora lub sprzętu warto
budować niezależne zestawy benchmarków. Mechanizm szybszy na
maszynie deweloperskiej może okazać się mniej korzystny na systemie
o innym cache, szerokości SIMD albo liczbie rdzeni.


\section*{Podsumowanie rozdziału}
\addcontentsline{toc}{section}{Podsumowanie rozdziału}

C++23 wyraźnie ulepsza projektowanie interfejsów (\code{expected}),
pracę na danych (\code{mdspan}), korutyny (\code{generator}) oraz
programowanie zakresowe. C++26 rozwija kontenery, wektoryzację,
wykonanie asynchroniczne i mechanizmy metaprogramowania.
Najważniejszą umiejętnością nie jest jednak mechaniczne używanie
najświeższej składni, lecz świadome odróżnianie gwarancji
standardu, stanu wdrożenia i wymagań aplikacji.

\section*{Literatura i dokumentacja}
\addcontentsline{toc}{section}{Literatura i dokumentacja}
\begingroup\sloppy
\begin{itemize}
\item \url{https://en.cppreference.com/w/cpp/23} --- zestawienie C++23.
\item \url{https://en.cppreference.com/w/cpp/26} --- zestawienie C++26.
\item \url{https://en.cppreference.com/w/cpp/compiler_support} --- wsparcie kompilatorów i bibliotek.
\item \url{https://en.cppreference.com/w/cpp/feature_test} --- makra do testowania funkcjonalności.
\item \url{https://en.cppreference.com/w/cpp/utility/expected} --- \code{std::expected}.
\item \url{https://en.cppreference.com/w/cpp/container/mdspan} --- \code{std::mdspan}.
\item \url{https://wg21.link} --- oficjalne dokumenty WG21 i propozycje standardu.
\end{itemize}
\endgroup
